\section{Complete even-characteristic symplectic cores in ranks two and three}
\label{sec:even-symplectic-low-rank}

\paragraph{Status and scope.}
The author has reviewed the complete derivation and executed all exact
certificates.  A separate whole-proof and primary-source peer review
has passed, together with an independent reconstruction of all
83 generic rows, 4,549 inequalities, 104 small-table rows and
3,442 spectra.  This section uses every ordinary irreducible character
in the two indicated symplectic ranks.  The complete official generic
tables are explicit mathematical inputs, including the correction
history recorded by their authors.  No claim about higher symplectic
rank, odd characteristic, or a partial unipotent table is made.

We retain $\nu(T)=a-\max_\lambda\dim\ker(T-\lambda)$ for an
$a$-dimensional finite-order complex operator.

\subsection*{D-ESC-01. Two complete generic character bounds}

\begin{lemma}[D-ESC-01: complete generic bounds]
\label{D-ESC-01}
For every nontrivial ordinary irreducible character of the indicated
group and every nonidentity element,
\[
\begin{array}{c|c|c}
\text{group and range}&|\chi(g)|/\chi(1)&\min\chi(1)\\\hline
\operatorname{Sp}_4(q),\ q\ge8\text{ even}&\le1/4&\ge196\\
\operatorname{Sp}_6(q),\ q\ge4\text{ even}&\le17/63&\ge378.
\end{array}                                                   \tag{ESC.1}
\]
Consequently every nonscalar finite-order linear normalizer of
either represented core has
\[
                           \nu(T)\ge23a/126.              \tag{ESC.2}
\]
\end{lemma}

\begin{proof}
The first complete table originates in Enomoto
\cite[\S4, pp.~80--85, and Tables~IV-1/IV-2,
pp.~92--94]{ESC:Enomoto}.  We use the official complete
\texttt{Sp4.0} implementation, not an unchecked transcription
of the original printed entries.  The official documentation
explicitly states that the printed paper has typographical
errors and that F.~L\"ubeck recomputed the implemented table
using Deligne--Lusztig theory \cite{G2T:GCT}.
That correction history is part of the input's provenance.

The second complete table is L\"ubeck's \texttt{Sp6.0}.
His full author text explicitly states the complete ordinary
character-table result, distinguishes it from the unipotent
subtable, and explains that the complete values were supplied
as CHEVIE computer files rather than printed in full
\cite[Introduction, pp.~3--4; \S2.2, pp.~18--21;
\S9, pp.~78--79]{ESC:Luebeck}.
The construction uses all Lusztig series and the additional
nonuniform class functions, so its scope is not merely the
first twelve unipotent character types.  The official
implementation is the entire \texttt{Sp6.0} table in
\cite{G2T:GCT}, at the same frozen commit as D-G2T-01.

Here is a finite exact proof of the new bounds from these
complete mathematical inputs.  As in D-G2T-01, write each
original entry as $\sum_\alpha p_\alpha(q)u_\alpha$ with
$p_\alpha\in\mathbf Q[q]$ and $|u_\alpha|=1$.
For each table put $x=q-q_0\ge0$ and write
\[
 p_\alpha(q_0+x)=\sum_k c_{\alpha k}x^k,
 \qquad d(q_0+x)=\sum_k d_kx^k.
\]
The same coefficientwise triangle inequality proves
$|\chi(g)|\le c\chi(1)$ as soon as
\[
                 c d_k-\sum_\alpha|c_{\alpha k}|\ge0
                      \quad\hbox{for every }k.             \tag{ESC.3}
\]
All of (ESC.3) pass with the following exact input parameters:
\[
\begin{array}{c|r|c|r|r|r}
\text{table}&q_0&c&\text{character/class types}&
\text{nontrivial cells}&\min d(q_0)\\\hline
\texttt{Sp4.0}&8&1/4&19&324&196\\
\texttt{Sp6.0}&4&17/63&66&4225&378.
\end{array}                                                   \tag{ESC.4}
\]
Every translated degree polynomial has nonnegative coefficients.
Thus its constant term gives the uniform degree lower bound.
Every nonidentity class-length polynomial likewise has
nonnegative translated coefficients and constant term greater
than one.  The trivial row and identity column are checked
against the separately declared degrees and class lengths.

The program \texttt{build\_even\_symplectic.py} checks the two
frozen source hashes and extracts the original mathematical
entry strings.  It executes no downloaded Julia code.
The offline original input has $436290$ bytes and SHA-256
\texttt{49ecbe27ad97b8cfe0204f96ef463ef1e8e0626ce6987bcccb0d8b33041e4fa0}.
The safe arithmetic parser already proved and reviewed in
D-G2T-01 constructs all $83$ nontrivial character-type rows
and $4549$ inequalities.  The deterministic coefficient
certificate has SHA-256
\texttt{92e32ae7b01d49469f31464472ae044b320b59a6400fedaec5003fde2760a3b3}.
\texttt{check\_even\_symplectic\_generic.py} reparses the
retained original input, matches every coefficient list, and
rechecks (ESC.3) using exact rational arithmetic.  All checks
pass.  The estimates are uniform in every admissible class
and character parameter, since only the absolute value one
of each phase is used.

For even $q$ both generic groups have trivial center, and the
published complete exceptional-multiplier list contains no
exception in these ranges \cite[Chapter~5, Table~1,
numbered p.~35]{BI:HHM}.  Thus a nontrivial irreducible
representation is faithful on its simple core.  Applying
the complex commutator estimate D-BI-18 to $c=17/63$ gives
\[
                       (1-c)/4=23/126.
\]
The first row has the stronger $3/16$ estimate, which is at
least $23/126$.  The commutator proof includes all normalizing
outer automorphisms and nonidentity scalar commutators, so
no outer-table assumption enters (ESC.2).
\end{proof}

\subsection*{D-ESC-02. The complete small-field endpoints}

\begin{lemma}[D-ESC-02: exact small-field normalizers]
\label{D-ESC-02}
In every nontrivial irreducible representation of
$\operatorname{Sp}_4(4)$, every nonscalar finite-order linear
normalizer satisfies $\nu(T)\ge a/3$.  Every such degree is
at least eighteen.

For every nonlinear irreducible representation of the universal
cover $2.\operatorname{Sp}_6(2)$, including those factoring
through its center, the following statements hold for every
nonscalar normalizer:
\begin{enumerate}
\item the unique degree-seven row has $\nu(T)\ge a/7$ and
has a real form with invariant symmetric bilinear form;
\item every other row has $\nu(T)\ge2a/9$;
\item among the nonlinear rows of degree at most nine, the
only degrees are seven and eight; the degree-eight row has
the stronger $\nu(T)\ge a/4$.
\end{enumerate}
\end{lemma}

\begin{proof}
We retain these four complete published ATLAS tables in CTblLib:
\[
\begin{array}{c|r|r|r}
\text{table}&\text{classes}&\text{selected rows}&
\text{nonscalar spectra}\\\hline
\texttt{S4(4)}&27&26&676\\
\texttt{S4(4).2}&30&20&580\\
\texttt{S4(4).4}&30&16&464\\
\texttt{2.S6(2)}&43&42&1722.
\end{array}                                                   \tag{ESC.5}
\]
The complete original input contains every character, class
size/order, prime power map, Frobenius--Schur indicator, and
actual fusion from the core to each extension.  Its size is
$19056$ bytes and SHA-256 is
\texttt{b51fae6eaa904bb5a1bdcac2f6aa85a940c9968ddba912bf115faa750e772fe2}.
The metadata record trivial multiplier and cyclic outer group
of order four for $\operatorname{Sp}_4(4)$, and multiplier
of order two with trivial outer group for
$\operatorname{Sp}_6(2)$.  The complete exceptional-multiplier
list agrees \cite[Chapter~5, Table~1, p.~35]{BI:HHM}.
The \texttt{S4(4).4} record additionally identifies its
construction as the full automorphism group.  These are
group-theoretic mathematical inputs in the standard CTblLib
metadata convention, not conclusions from numerical
orthogonality alone.

The author GAP exporter checks full row orthogonality and
every cyclic spectrum.  The separate polynomial Fourier
verifier \texttt{check\_even\_symplectic\_small.py} selects
irreducible extension restrictions anew by exact character
equality, verifies both normal fusions and cyclic quotients,
and reconstructs every eigenvalue multiplicity using (G2T.8).
Every multiplicity is a nonnegative integer and their sum
is the degree.  All $104$ selected rows and $3442$ spectra
pass.  The three $\operatorname{Sp}_4(4)$ tables have minimum
ratio $1/3$ and minimum degree eighteen.

In the universal $\operatorname{Sp}_6(2)$ table the only
quarter-bound failures are exactly
\[
\begin{array}{r|r|c}
\text{row}&a&\min\nu/a\\\hline
2&7&1/7\\
5&21&5/21\\
6&27&2/9.
\end{array}                                                   \tag{ESC.6}
\]
The other rows satisfy the quarter bound, so all non-seven
rows satisfy $2/9$.  The degree-seven character takes rational
integer values; the verifier recomputes its indicator directly:
\[
             \frac1{|U|}\sum_{g\in U}\chi(g^2)=1,
                     \qquad U=2.\operatorname{Sp}_6(2).
\]
The standard Frobenius--Schur criterion supplies its invariant
symmetric bilinear form.  The complete degree list has exactly
the two asserted nonlinear degrees at most nine.

For $\operatorname{Sp}_4(4)$ the cyclic subgroups of the outer
group have orders one, two and four, all represented in
(ESC.5).  Given an actual normalizer, adjust its automorphism
by an inner action and choose a corresponding element $t$
of the appropriate fixed extension $A$.  If $A/U$ has order
$d$, Schur's lemma gives $T^d=c\rho(t^d)$.  Rescaling $T$
by a $d$th root of $c^{-1}$ gives an honest representation
of that fixed extension, irreducible on its core.  The
original operator is scalar times an element of a tested
row.  Thus no splitting or implicit choice of scalar power
relation is assumed.  For $\operatorname{Sp}_6(2)$ there is
no outer action.  Uniqueness of automorphism lifting to its
universal perfect cover implies that every normalizer is
scalar times a core element.  Retaining every nonlinear
universal-cover row includes every possible central kernel.
The stated finite bounds therefore hold for all actual
normalizers.
\end{proof}

\subsection*{D-ESC-03. Exact source consequences}

\begin{theorem}[D-ESC-03: characteristic-two rank-two/rank-three cores]
\label{D-ESC-03}
Let $J$ be the finite image of an inducing block in a finite-cover
Weil-sixfold source.
\begin{enumerate}
\item If $J$ has an irreducibly acting normal central product
of quasisimple factors with simple quotients
$\operatorname{Sp}_4(q)$ or $\operatorname{Sp}_6(q)$,
where all $q\ge4$ are even, its actual period rank is at
most eight.  Mixed ranks and field sizes are allowed.
\item Suppose instead that the simple quotients of all
factors are $\operatorname{Sp}_6(2)$ and $J$ permutes the
factors transitively.  If the source has period rank nine,
there is exactly one factor and the inducing degree is seven.
\item In the nonsplit problem over any of the eight
non-Eisenstein imaginary quadratic fields of class number
one, that last degree-seven case is excluded by D-BI-17.
Consequently no transitive normal core of either of these
two symplectic ranks in characteristic two remains in that
problem.
\end{enumerate}
All parts allow arbitrary scalar phases, compatible outer
automorphisms and tensor permutations, and a larger inducing
character field.  Parts (1)--(2) impose no nonsplit determinant
condition.  Part (2) is an exact remaining-case statement,
not an existence assertion for its endpoint.
\end{theorem}

\begin{proof}
For part (1), write the irreducible core representation as
$V=\bigotimes V_i$, with dimensions $a_i$, and put
$e=\prod a_i$.  The intrinsic simple factors and their full
matrix algebras are permuted by $J$.  If an operator preserves
all factors and is nonscalar, a nonscalar factor intertwiner
controls its largest eigenspace fraction.  D-ESC-01 and
D-ESC-02 give $\nu(T)\ge23e/126$.  The minimum factor degree
in the combined list is eighteen.

A nontrivial cycle of $l\ge2$ tensor factors of dimension
$a$ is conjugate to cyclic word rotation with a diagonal
twist at one position.  Its constant basis words number
$a$, and every other orbit has at least two words.  A fixed
eigenvalue occurs at most once on each weighted orbit, so
its eigenspace has dimension at most
$a+(a^l-a)/2\le3a^l/4$.  Tensoring with the other cycles
preserves this fraction.  Hence factor permutations give
the stronger $\nu(T)\ge e/4$.  Genuine nonidentity scalar
inertia has full moving rank $e$.

If the actual period rank were nine, the intermediate-cover
factorization and Euler argument of D-BI-01 give genus zero,
at least twelve genuine inertias, and total moving rank
$2e+6$.  Indeed genus one has at most six marked-base
parameters and genus at least two would make
$6+2e(1-h)$ negative.  Therefore
\[
                  12\frac{23e}{126}\le2e+6,
                        \qquad e\le63/2.                 \tag{ESC.7}
\]
For at least two factors, $e\ge18^2>63/2$.
For a single generic factor the minimum degrees $196,378$
also violate (ESC.7).  A single $\operatorname{Sp}_4(4)$
factor has the stronger one-third bound, giving
$4e\le2e+6$, or $e\le3<18$.  This proves part (1).

For part (2), transitivity makes all factor dimensions equal,
say $a$, and $e=a^t$.  First take $a\ne7$.
The factor-preserving bound is $2e/9$ and the factor-permuting
bound is $e/4\ge2e/9$.  The same actual branch budget gives
\[
                        12\frac{2e}{9}\le2e+6,
                              \qquad e\le9.               \tag{ESC.8}
\]
The full degree list in D-ESC-02 forces $t=1$ and $a=8$.
Its stronger quarter bound instead gives
$12\lceil8/4\rceil=24>22=2e+6$, a contradiction.

Now let $a=7$.  Every factor has an invariant symmetric
bilinear form, so their tensor product is symmetric and
spans the unique invariant bilinear-form line of the
irreducible normal core.  The full normalizer preserves
that line, with a similitude character $\mu$.
For the original quadratic field $K$ and the inducing
character field $L$, induction gives $K\subset L$.
If $\mu$ were trivial, every finite-order matrix would
preserve a symmetric bilinear form exactly; its eigenvalues
would pair with their inverses, and its trace would be real.
This would contradict the imaginary quadratic subfield of
$L$.  Hence $\mu$ is nontrivial.

The genuine inertias generate and have product one, so at
least two have nontrivial multiplier.  Their fixed spaces
are totally isotropic, and each moves at least $e/2$
dimensions.  All other genuine inertias move at least
$e/7$ by D-ESC-02 and the tensor-cycle bound.  Thus
\[
                    2e+6\ge e+10e/7,
                              \qquad e\le14.              \tag{ESC.9}
\]
Since $e=7^t$, this forces $t=1$ and $e=7$, proving part (2).
Part (3) is now precisely the previously proved D-BI-17
exclusion of every degree-seven inducing block in its
stated eight-field nonsplit scope.  It does not require
identifying a new polarization or assuming equality $K=L$.
\end{proof}

\paragraph{Boundary and unsuccessful stronger estimates.}
The nonsimple group $\operatorname{Sp}_4(2)$ is not included;
its derived simple group is $A_6$ and belongs to D-BI-21.
The initial attempted quarter-character bounds at threshold
$q=4$ failed at six cells for \texttt{Sp4.0} and three cells
for \texttt{Sp6.0}.  Those failures are retained in
\texttt{even\_symplectic\_exploration.json}.  The proved
statement uses the different threshold and exact ratio in
(ESC.4), together with complete small-group spectra.
The degree-seven exception is retained in part (2), and
is removed only under the precise earlier theorem's scope.
